\begin{table}[t]
\centering
\caption{Summary Statistics: Dissolved Oxygen by TMDL Coverage}
\label{tab:summary}
\begin{threeparttable}
\begin{tabular}{lcccc}
\toprule
 & \multicolumn{2}{c}{High TMDL Coverage} & \multicolumn{2}{c}{Low TMDL Coverage} \\
\cmidrule(lr){2-3} \cmidrule(lr){4-5}
 & Mean DO & N & Mean DO & N \\
\midrule
Pre-TMDL (2000--2009) & 9.25 & 8,013 & 7.65 & 8,085 \\
Post-TMDL (2010--2023) & 9.14 & 7,171 & 7.93 & 7,314 \\
\midrule
Stations & \multicolumn{2}{c}{3,335} & \multicolumn{2}{c}{2,758} \\
HUC8 Watersheds & \multicolumn{2}{c}{18} & \multicolumn{2}{c}{24} \\
Mean TMDL Share & \multicolumn{2}{c}{0.48} & \multicolumn{2}{c}{0.04} \\
\bottomrule
\end{tabular}
\begin{tablenotes}
\item \textit{Notes:} Summary statistics for dissolved oxygen (DO) readings from the EPA Water Quality Portal, 2000--2023. Stations are matched to HUC8 watersheds and classified by TMDL coverage using EPA ATTAINS data. High TMDL Coverage denotes stations in HUC8 watersheds where the share of impaired assessment units with completed TMDLs exceeds the sample median. DO measured in mg/L.
\end{tablenotes}
\end{threeparttable}
\end{table}

